\documentclass[11pt,a4paper]{article}
\usepackage[utf8]{vietnam}
\usepackage{amsmath,amssymb}
\usepackage{tabularx}
\usepackage{array}
\usepackage{booktabs}
\usepackage{multirow}
\usepackage[landscape, top=1.2cm, bottom=1.2cm, left=1.2cm, right=1.2cm]{geometry}
\usepackage{fancyhdr}

\pagestyle{fancy}
\fancyhf{}
\rfoot{\small\textit{Trang \thepage}}
\lfoot{\small\textit{Ma trận và Bản đặc tả Đề kiểm tra Giữa HK I -- Toán 7 -- Năm học 2026--2027}}
\renewcommand{\headrulewidth}{0pt}

\newcolumntype{Y}{>{\centering\arraybackslash}X}
\newcolumntype{M}[1]{>{\centering\arraybackslash}p{#1}}
\newcolumntype{L}[1]{>{\raggedright\arraybackslash}p{#1}}

\begin{document}

\begin{center}
  \textbf{\large MA TRẬN, BẢN ĐẶC TẢ ĐỀ KIỂM TRA GIỮA HỌC KÌ I}\\[2mm]
  \textbf{\large MÔN TOÁN -- LỚP 7 -- NĂM HỌC 2026 -- 2027}\\[1mm]
  \textit{(Căn cứ theo Công văn 7991/BGDĐT-GDTrH -- Áp dụng chung cho Đề 1 và Đề 2)}
\end{center}

\vspace{1.5mm}
\noindent\textbf{\large I. KHUNG MA TRẬN ĐỀ KIỂM TRA (Bước 1)}
\vspace{1mm}

\noindent
\small
\begin{tabularx}{\linewidth}{|M{7mm}|L{2.8cm}|X|*{12}{M{7mm}|}M{1.1cm}|M{1.3cm}|}
\hline
\multirow{3}{*}{\textbf{TT}} & \multirow{3}{*}{\textbf{Chủ đề / Chương}} & \multirow{3}{*}{\textbf{Nội dung / đơn vị kiến thức}} & \multicolumn{12}{c|}{\textbf{Mức độ đánh giá}} & \multicolumn{2}{c|}{\textbf{Tổng}} \\ \cline{4-17}
& & & \multicolumn{3}{c|}{\textbf{Nhiều lựa chọn}} & \multicolumn{3}{c|}{\textbf{Đúng -- Sai (ý)}} & \multicolumn{3}{c|}{\textbf{Trả lời ngắn}} & \multicolumn{3}{c|}{\textbf{Tự luận}} & \textbf{Số} & \textbf{Điểm} \\ \cline{4-15}
& & & \textbf{B} & \textbf{H} & \textbf{VD} & \textbf{B} & \textbf{H} & \textbf{VD} & \textbf{B} & \textbf{H} & \textbf{VD} & \textbf{B} & \textbf{H} & \textbf{VD} & \textbf{câu/ý} & \textbf{(\%)} \\ \hline

1 & \textbf{Chương I: Số hữu tỉ} & Tập hợp số hữu tỉ. Thứ tự trong $\mathbb{Q}$ (Bài 1) & 2 & & & & & & & & & & & & 2 & 0,5 (5\%) \\ \cline{3-17}
& & Các phép tính với số hữu tỉ (Bài 2, 3, 4) & 2 & & & & & & & 1 & & & 1 & 2 & 6 & 3,0 (30\%) \\ \hline

2 & \textbf{Chương II: Số thực} & Số thập phân vô hạn tuần hoàn. Làm tròn (Bài 5) & 1 & 1 & & & & & & & & & & & 2 & 0,5 (5\%) \\ \cline{3-17}
& & Số vô tỉ. Căn bậc hai số học. Số thực (Bài 6, 7) & 2 & & & 1 & 2 & 1 & 1 & & 1 & & & & 8 & 2,5 (25\%) \\ \hline

3 & \textbf{Chương III: Góc \& ĐT song song} & Góc ở vị trí đặc biệt. Tia phân giác (Bài 8, 9) & 1 & 1 & & & & & & & & & & & 2 & 0,5 (5\%) \\ \cline{3-17}
& & Hai đường thẳng song song (Bài 10) & 2 & & & 2 & 1 & 1 & & 1 & & & 1 & & 8 & 2,5 (25\%) \\ \cline{3-17}
& & Định lí và chứng minh định lí (Bài 11) & & & & & & & & & & & & 1 & 1 & 0,5 (5\%) \\ \hline

\multicolumn{3}{|c|}{\textbf{Tổng số câu / ý}} & \textbf{8} & \textbf{2} & \textbf{0} & \textbf{3} & \textbf{3} & \textbf{2} & \textbf{1} & \textbf{2} & \textbf{1} & \textbf{0} & \textbf{2} & \textbf{3} & \textbf{29} & \textbf{10,0} \\ \hline
\multicolumn{3}{|c|}{\textbf{Tổng số điểm}} & \multicolumn{3}{c|}{\textbf{3,0 điểm}} & \multicolumn{3}{c|}{\textbf{2,0 điểm}} & \multicolumn{3}{c|}{\textbf{2,0 điểm}} & \multicolumn{3}{c|}{\textbf{3,0 điểm}} & \multicolumn{2}{c|}{\textbf{10,0 đ}} \\ \hline
\multicolumn{3}{|c|}{\textbf{Tỉ lệ \%}} & \multicolumn{3}{c|}{\textbf{30\%}} & \multicolumn{3}{c|}{\textbf{20\%}} & \multicolumn{3}{c|}{\textbf{20\%}} & \multicolumn{3}{c|}{\textbf{30\%}} & \multicolumn{2}{c|}{\textbf{100\%}} \\ \hline
\end{tabularx}

\vspace{2.5mm}
\noindent\textit{* Ghi chú tỉ lệ nhận thức:} Nhận biết: 4,0 điểm (40\%) -- Thông hiểu: 3,0 điểm (30\%) -- Vận dụng: 3,0 điểm (30\%).\\
\textit{* Phân môn:} Đại số 6,5 điểm (65\%) -- Hình học 3,5 điểm (35\%).

\newpage
\noindent\textbf{\large II. BẢN ĐẶC TẢ ĐỀ KIỂM TRA GIỮA HỌC KÌ I (Bước 2)}
\vspace{2mm}

\noindent
\small
\begin{tabularx}{\linewidth}{|M{7mm}|L{2.8cm}|L{3.2cm}|M{1.4cm}|X|L{3.2cm}|M{1.1cm}|M{1.6cm}|}
\hline
\textbf{TT} & \textbf{Chương/Chủ đề} & \textbf{Nội dung} & \textbf{Mức độ} & \textbf{Yêu cầu cần đạt} & \textbf{Dạng câu / Câu số} & \textbf{Số câu} & \textbf{Năng lực} \\ \hline

1 & \textbf{Số hữu tỉ} (Chương I) & Tập hợp số hữu tỉ. Thứ tự trong $\mathbb{Q}$ & Biết & Nhận biết số hữu tỉ, kí hiệu $\mathbb{Q}$, số đối, thứ tự trong $\mathbb{Q}$. & NLC: C1, C2 & 2 & TD, GT \\ \cline{3-8}
& & Phép tính với số hữu tỉ & Biết & Luỹ thừa với số mũ tự nhiên; quy tắc chuyển vế, dấu ngoặc. & NLC: C3, C4 & 2 & TD \\ \cline{4-8}
& & & Hiểu & Tính toán thành thạo, tìm $x$, phối hợp các phép tính. & TLN: C16;\; TL: 19a & 2 & TD, CC \\ \cline{4-8}
& & & VD & Vận dụng giải bài toán tìm $x$ luỹ thừa, bài toán thực tế tỉ số/phần trăm. & TL: 19b, C20 & 2 & GQ, CC \\ \hline

2 & \textbf{Số thực} (Chương II) & Số thập phân VHTH. Làm tròn số & Biết & Nhận biết chu kì của số thập phân vô hạn tuần hoàn. & NLC: C5 & 1 & TD \\ \cline{4-8}
& & & Hiểu & Ước lượng và làm tròn số theo độ chính xác cho trước. & NLC: C8 & 1 & TD, CC \\ \cline{3-8}
& & Căn bậc hai số học. Số thực & Biết & Căn bậc hai số học, số vô tỉ, giá trị tuyệt đối. & NLC: C6, C7;\; ĐS: 13a;\; TLN: C15 & 4 & TD \\ \cline{4-8}
& & & Hiểu & So sánh các số thực, tính chất số thập phân vô hạn tuần hoàn. & ĐS: 13b, 13c & 2 & TD, CC \\ \cline{4-8}
& & & VD & Giải phương trình chứa dấu giá trị tuyệt đối; bài toán thực tế căn bậc hai. & ĐS: 13d;\; TLN: C18 & 2 & TD, GQ \\ \hline

3 & \textbf{Góc \& ĐT song song} (Chương III) & Góc đặc biệt. Tia phân giác & Biết & Tính chất hai góc đối đỉnh; khái niệm tia phân giác. & NLC: C9 & 1 & TD \\ \cline{4-8}
& & & Hiểu & Tính số đo góc kề bù, góc đối đỉnh, góc tạo bởi tia phân giác. & NLC: C11 & 1 & TD, GT \\ \cline{3-8}
& & Hai ĐT song song & Biết & Tiên đề Euclid; nhận biết góc so le trong, đồng vị; quan hệ từ vuông góc đến song song. & NLC: C10, C12;\; ĐS: 14a, 14b & 4 & TD, GT \\ \cline{4-8}
& & & Hiểu & Dấu hiệu nhận biết hai đường thẳng song song; tính số đo góc tạo bởi cát tuyến. & ĐS: 14c;\; TLN: C17;\; TL: 21a & 3 & TD, GT \\ \cline{4-8}
& & & VD & Vận dụng tia phân giác và song song tính góc phức hợp. & ĐS: 14d & 1 & TD, GQ \\ \cline{3-8}
& & Định lí \& Chứng minh định lí & VD & Chứng minh bài toán hình học có lập luận chặt chẽ (song song, vuông góc). & TL: 21b & 1 & TD, GQ \\ \hline
\end{tabularx}

\newpage
\noindent\textbf{\large III. BƯỚC 5: BẢNG KIỂM TRA LẠI ĐỀ (Áp dụng chung cho Đề 1 và Đề 2)}
\vspace{2mm}

\noindent
\normalsize
\begin{tabularx}{\linewidth}{|M{1.2cm}|L{4.2cm}|X|M{2.4cm}|}
\hline
\textbf{STT} & \textbf{Tiêu chí kiểm tra} & \textbf{Kết quả đối soát thực tế} & \textbf{Kết luận} \\ \hline
1 & \textbf{Độ cân đối giữa các mức độ nhận thức} & Biết 4,0 đ (40\%) -- Hiểu 3,0 đ (30\%) -- Vận dụng 3,0 đ (30\%). Đúng chuẩn quy định tại Công văn 7991. Dạng câu: Nhiều lựa chọn 3,0 đ; Đúng/Sai 2,0 đ; Trả lời ngắn 2,0 đ; Tự luận 3,0 đ. & \textbf{ĐẠT} \\ \hline
2 & \textbf{Tỉ lệ phân môn} & Đại số 6,5 đ (65\%) -- Hình học 3,5 đ (35\%). Phù hợp tiến độ dạy học đến tuần 10 (Đại số 24 tiết $\approx 69\%$, Hình học 11 tiết $\approx 31\%$). & \textbf{ĐẠT} \\ \hline
3 & \textbf{Sự phù hợp chuẩn đầu ra} & Mọi câu hỏi bám sát chuẩn kiến thức, kĩ năng môn Toán 7 (CT GDPT 2018). Không sử dụng các kiến thức chưa học (tam giác bằng nhau, định lí Pythagore). & \textbf{ĐẠT} \\ \hline
4 & \textbf{Tính chính xác của đáp án} & Đã giải độc lập cả hai mã đề. Các câu trắc nghiệm chỉ có 1 phương án đúng duy nhất; lời giải tự luận phân bậc $0{,}25$đ rõ ràng. & \textbf{ĐẠT} \\ \hline
\end{tabularx}

\end{document}
